\section{Дані}
\label{sec:data}

\subsection{Вихідний корпус}

ЄДРСР публікує всі рішення українських судів починаючи з 2005 року.
Ми обробляємо 101 мільйон повнотекстових рішень, розподілених за роком ухвалення (2007–2026).
Видобування цитувань використовує скомпільовані регулярні вирази для 13 кодексів із нормалізацією скорочень та розгортанням діапазонів статей.
Первинне видобування дає 502 мільйони записів цитувань усіх типів (статті кодексів, посилання на справи, конституційні цитування, посилання на закони за номером).
Після обмеження до цитувань статей кодексів та збереження лише коректно сформованих посилань на статті ми отримуємо 396 мільйонів валідованих записів – основу нашої оцінки.

\subsection{Річні зрізи}

Для кожного року $y \in \{2007, \ldots, 2026\}$ ми будуємо оціночний зріз:
\begin{enumerate}[nosep,leftmargin=*]
  \item Видобуваємо всі цитування кодексів із рішень, ухвалених у році $y$.
  \item Фільтруємо статті: мінімум 50 цитувань, обмеження до 5{,}000 найчастіших.
  \item Фільтруємо справи: від 3 до 200 процитованих статей на рішення.
\end{enumerate}
Зріз 2024 року містить 3{,}671 статтю, 1{,}801{,}481 справу та 16,4M ребер цитування.

\section{Метод}
\label{sec:method}

\subsection{Протокол оцінювання}

Нехай $M \in \{0,1\}^{|U| \times |V|}$ – бінарна матриця інцидентності справа–стаття, а $C = M^\top M$ (з $\text{diag}(C) = 0$) – матриця спільного цитування статей.

Для кожної справи $u$ з процитованими статтями $S_u$ ($|S_u| \geq 3$) та кожної цільової $v_t \in S_u$:
оцінюємо всіх кандидатів $v \notin S_u \cup \{v_t\}$ та обчислюємо ранг $v_t$.
Оскільки $\text{diag}(C) = 0$, власний внесок цільової статті природно виключається.

Ми відбираємо 200{,}000 справ на рік, що дає $\sim$1,8M передбачень на зріз.

\subsection{Функції оцінювання}

\textbf{Спільні сусіди}: $\text{CN}(v_t \mid S) = \sum_{s \in S} C_{s, v_t}$

\textbf{Адамік–Адар}: $\text{AA}(v_t \mid S) = \sum_{s \in S} C_{s, v_t} / \max(\log d_s, 1)$

\textbf{Базова лінія за ступенем}: $\text{DG}(v_t) = d_{v_t}$ (лише популярність).

\subsection{Контрольні абляції}

Для відокремлення справжнього темпорального згасання від конфаундерів ми вводимо два контрольні експерименти:

\paragraph{Абляція з фіксованим набором статей.}
Ми визначаємо статті з $\geq$50 цитуваннями \emph{як} у 2012 (пік), \emph{так і} в 2024 (нещодавній період), отримуючи спільний словник.
Оцінювання лише на цьому фіксованому наборі ізолює темпоральне згасання від композиційного зсуву (нові статті, що розмивають продуктивність).

\paragraph{Темпоральне розділення навчання/тестування.}
У межах кожного року ми будуємо $C$ за першими 50\% справ (за ідентифікатором документа) та оцінюємо на других 50\%.
Це усуває витік даних: матриця спільного цитування ніколи не бачить оціночних справ.
